% =====================================================================
% TEMPLATE ELABORASI SATU KOMPONEN
% Salin file ini menjadi komponen_NAMA.tex, lalu ganti placeholder.
% Struktur ini diselaraskan secara generik mengikuti alur dht22.tex.
% Jangan meng-input file template ini secara langsung ke bab.
% =====================================================================

\subsection{Nama Komponen -- Peran/Kategori Singkat}
\label{subsec:nama_komponen}

% ---------------------------------------------------------------------
\subsubsection{Identitas dan Fungsi}
Jelaskan modul/komponen secara ringkas, fungsinya dalam sistem tertanam,
kategori komponen (sensor analog/digital, aktuator, modul komunikasi, atau IC driver),
serta alasan komponen dipilih (misal: akurasi tinggi, antarmuka sederhana, konsumsi daya rendah,
ketersediaan pustaka, atau rentang kerja yang lebar).

\begin{table}[H]
	\centering
	\caption{Identitas komponen Nama Komponen}
	\label{tab:nama_komponen_identitas}
	\small
	\begin{tabularx}{\textwidth}{@{}lX@{}}
		\toprule
		Parameter          & Keterangan                                                       \\
		\midrule
		Nama komponen      & Nama dagang / nama populer komponen                              \\
		Model/part number  & Part number resmi pabrikan (misal: AM2302, HC-SR04, SG90, DS18B20) \\
		Jenis              & Sensor (Analog/Digital) / Aktuator / Modul Antarmuka             \\
		Fungsi utama       & Deskripsi singkat fungsi utama komponen                          \\
		Sumber dokumentasi & Judul/tautan datasheet atau referensi pabrikan resmi             \\
		\bottomrule
	\end{tabularx}
\end{table}

% ---------------------------------------------------------------------
\subsubsection{Prinsip Kerja}
Jelaskan fenomena fisika atau prinsip transduser elektronika yang memungkinkan komponen
beroperasi. Misalnya: efek termoelektrik/piezoelektrik, perubahan resistansi atau kapasitansi,
efek medan magnet Hall, fotodioda optik, atau induksi elektromagnetik pada motor/aktuator.
Uraikan elemen-elemen penyusun internalnya:

\begin{enumerate}
	\item \textbf{Elemen Transduser Utama:}
	      Jelaskan transduser primer yang berinteraksi langsung dengan besaran fisik luar.
	      Uraikan bagaimana perubahan besaran lingkungan memengaruhi karakteristik fisis/listrik elemen ini.

	      \begin{figure}[H]
		      \centering
		      % Ganti dengan path file kurva/karakteristik transduser (misal kurva respons, diagram blok internal)
		      % \includegraphics[width=0.50\textwidth]{images/karakteristik_nama_komponen.png}
		      \fbox{\parbox[c][3.5cm][c]{0.50\textwidth}{\centering Tempatkan kurva karakteristik atau diagram blok elemen di sini}}
		      \caption{Kurva karakteristik atau struktur internal elemen komponen Nama Komponen.}
		      \label{fig:nama_komponen_karakteristik}
	      \end{figure}

	\item \textbf{Elemen Pengkondisi Sinyal / Sirkuit Internal:}
	      Jelaskan pengkondisian sinyal internal seperti penguat (op-amp), filter, osilator,
	      ADC internal, atau mikrokontroler internal yang memproses besaran dari transduser primer.
\end{enumerate}

\paragraph{Konsep dan Formulasi Matematis Terkait}
Jika terdapat hukum fisika, hubungan transfer function, atau konversi matematis besaran,
tuliskan persamaannya secara eksplisit:

\begin{equation}
	y = f(x)
	\label{eq:nama_komponen_formulasi}
\end{equation}

\noindent di mana:
\begin{itemize}
	\item $y$ = Variabel keluaran (misal tegangan keluaran, besaran digital terhitung, dsb.).
	\item $x$ = Variabel masukan fisik (misal suhu, jarak, sudut, intensitas cahaya, dsb.).
\end{itemize}

% ---------------------------------------------------------------------
\subsubsection{Alur Kerja Internal dan Protokol Komunikasi}
Alur pemrosesan data dan interaksi komponen mengikuti pola \textbf{Input $\rightarrow$ Proses $\rightarrow$ Output}:
\begin{enumerate}
	\item \textbf{Input:} Stimulus besaran fisik dari lingkungan (untuk sensor) atau pulsa perintah/sinyal kontrol dari mikrokontroler (untuk aktuator).
	\item \textbf{Proses:} Pemrosesan internal oleh sirkuit terpadu komponen (penguatan sinyal, konversi ADC, kompensasi kalibrasi OTP/EEPROM, atau modulasi PWM driver daya).
	\item \textbf{Output:} Sinyal listrik terformat (tegangan analog, pulsa logika digital, paket serial I$^2$C/SPI/UART/1-Wire) atau aksi mekanik/fisik langsung.
\end{enumerate}

\paragraph{Detail Sinyal Protokol Komunikasi}
Jelaskan urutan timing, sinyal jabat tangan (\textit{handshake}), atau format transmisi data:
\begin{itemize}
	\item \textbf{Sinyal Start / Inisiasi:} Kondisi bus atau pulsa pemicu yang memulai transmisi data.
	\item \textbf{Sinyal Respons / Acknowledgment:} Respons dari komponen yang mengonfirmasi kesiapan komunikasi.
	\item \textbf{Format Transmisi Data:} Struktur data per bit/byte, durasi pulsa logika LOW/HIGH, atau timing clock.
	\item \textbf{Struktur Frame / Pemeriksaan Galat:} Pembagian paket data (misal Data MSB/LSB + Checksum/CRC jika ada).
\end{itemize}

% Diagram Pewaktu / Timing Diagram (opsional menggunakan TikZ atau gambar ilustrasi)
\begin{figure}[htbp]
	\centering
	\begin{tikzpicture}[x=0.18mm, y=10mm, >=Stealth, font=\sffamily\small]
		% Base axis / waktu
		\draw[->, color=gray!60, thick] (0, 0) -- (600, 0) node[right, color=black] {\textit{Waktu}};
		\draw[->, color=gray!60, thick] (0, 0) -- (0, 2.3) node[above, color=black] {Sinyal};

		% Level logika
		\node[left] at (0, 1.8) {\textbf{HIGH}};
		\node[left] at (0, 0.4) {\textbf{LOW}};
		\draw[dotted, color=gray!40] (0, 1.8) -- (580, 1.8);
		\draw[dotted, color=gray!40] (0, 0.4) -- (580, 0.4);

		% Bentuk gelombang sinyal ilustratif
		\draw[ultra thick, color={rgb,255:red,30;green,64;blue,175}]
		(0, 1.8) -- (50, 1.8) -- (50, 0.4) -- (180, 0.4) -- (180, 1.8) -- (240, 1.8)
		-- (240, 0.4) -- (330, 0.4) -- (330, 1.8) -- (400, 1.8) -- (400, 0.4)
		-- (450, 0.4) -- (450, 1.8) -- (550, 1.8);

		% Anotasi rentang waktu
		\draw[stealth-stealth, color={rgb,255:red,30;green,64;blue,175}] (50, 2.0) -- (180, 2.0) node[midway, above, font=\scriptsize] {Inisiasi / Trigger};
		\draw[stealth-stealth, color={rgb,255:red,15;green,118;blue,110}] (240, 2.0) -- (330, 2.0) node[midway, above, font=\scriptsize] {Respons};
		\draw[stealth-stealth, color={rgb,255:red,217;green,119;blue,6}] (400, 2.0) -- (550, 2.0) node[midway, above, font=\scriptsize] {Transmisi Data / Pulsa};
	\end{tikzpicture}
	\caption{Diagram pewaktu (\textit{timing diagram}) atau bentuk sinyal komponen Nama Komponen.}
	\label{fig:nama_komponen_timing}
\end{figure}

Pada implementasi praktis mikrokontroler, protokol tingkat rendah ini umumnya dienkapsulasi
melalui pustaka perangkat lunak (\textit{driver library}) yang menangani pewaktuan mikrodetik,
pengalamatan register bus, serta kalkulasi integritas data secara otomatis sehingga pengguna
cukup memanggil fungsi API tingkat tinggi.

% ---------------------------------------------------------------------
\subsubsection{Besaran yang Diukur atau Dikendalikan}
Jelaskan besaran fisik utama yang menjadi domain kerja komponen (misal suhu, kelembapan,
jarak, sudut mekanik, laju aliran, konsentrasi gas, dsb.):

\begin{table}[H]
	\centering
	\caption{Karakteristik besaran komponen Nama Komponen}
	\label{tab:nama_komponen_karakteristik_besaran}
	\small
	\begin{tabularx}{\textwidth}{@{}lXX@{}}
		\toprule
		Parameter                            & Nilai                   & Keterangan                           \\
		\midrule
		Besaran utama                        & Besaran fisik utama     & Parameter operasional target         \\
		Rentang kerja                        & Nilai min s.d. maks     & Rentang dinamis pengukuran / kendali \\
		Resolusi                             & Nilai per satuan        & Resolusi bit atau sensitivitas diskret \\
		Akurasi                              & Toleransi kesalahan     & Batas galat pengukuran/posisi        \\
		Pengulangan (\textit{Repeatability}) & Nilai toleransi ulang   & Presisi pada kondisi berulang        \\
		Stabilitas / Waktu respons           & Nilai durasi/stabilitas & Kecepatan tanggapan terhadap step input \\
		\bottomrule
	\end{tabularx}
\end{table}

% ---------------------------------------------------------------------
\subsubsection{Karakteristik Sinyal}
Uraikan format sinyal elektrik antarmuka (analog tegangan/arus, pulsa PWM, atau paket serial digital):

\begin{equation}
	\text{Sinyal / Paket} = \text{Frame Header} + \text{Payload Data} + \text{Checksum}
	\label{eq:nama_komponen_packet}
\end{equation}

\begin{table}[H]
	\centering
	\caption{Karakteristik sinyal Nama Komponen}
	\label{tab:nama_komponen_sinyal}
	\small
	\begin{tabularx}{\textwidth}{@{}lXX@{}}
		\toprule
		Aspek                       & Nilai                                    & Keterangan                                       \\
		\midrule
		Jenis sinyal                & Analog / Digital Serial / PWM / I$^2$C   & Format media transmisi sinyal                    \\
		Level logika                & TTL 0--5\,V atau 3,3\,V LVCMOS           & Level tegangan high/low                          \\
		Rentang tegangan operasional& 3,3--5,0\,V DC                           & Rentang toleransi tegangan sinyal                \\
		Frekuensi / Kecepatan Data  & Nilai frekuensi / baud rate / sampling   & Laju pembaruan atau frekuensi clock sinyal       \\
		Sinyal Inisiasi / Pemicu    & Durasi / kondisi inisiasi                & Perintah start dari MCU ke komponen              \\
		Respons Komponen            & Durasi / pola handshake                  & Konfirmasi kesiapan dari komponen                \\
		Format Data / Encoding      & Format bit 0/1 atau lebar pulsa          & Representasi data biner atau analog              \\
		\bottomrule
	\end{tabularx}
\end{table}

% ---------------------------------------------------------------------
\subsubsection{Pin, Catu Daya, dan Antarmuka}
Jelaskan varian kemasan fisik komponen (misal: IC \textit{standalone} mentah vs modul \textit{breakout board}
yang telah dilengkapi resistor pull-up, kapasitor filter, dan pin header).

\begin{table}[H]
	\centering
	\caption{Pemetaan pin dan antarmuka komponen Nama Komponen}
	\label{tab:nama_komponen_pin}
	\small
	\begin{tabularx}{\textwidth}{@{}lllXX@{}}
		\toprule
		Pin (IC) & Pin (Modul Breakout) & Kabel   & Fungsi   & Keterangan                              \\
		\midrule
		Pin 1    & VCC / VIN            & Merah   & $V_{CC}$ & Catu daya positif (3,3\,V / 5\,V)       \\
		Pin 2    & OUT / DATA / SDA     & Kuning  & Sinyal   & Jalur sinyal analog/digital bi-direksional \\
		Pin 3    & SCL / TRIG / NC      & Putih   & Kontrol  & Pin clock / pemicu / \textit{Not Connected} \\
		Pin 4    & GND                  & Hitam   & GND      & Referensi ground (0\,V)                 \\
		\bottomrule
	\end{tabularx}
\end{table}

Jelaskan bila terdapat pin konfigurasi alamat (\textit{address select}), pin interupsi,
atau resistor pull-up/down bawaan pada modul \textit{breakout}.

\subsubsection{Kebutuhan Catu Daya}
Tegangan kerja optimal komponen berada pada rentang $3{,}3\text{--}5{,}0\,\text{V DC}$.
Konsumsi arus berada pada kisaran nominal saat aktif dan mode siaga (\textit{standby/sleep}).
Untuk menjaga kestabilan sinyal dan mencegah fluktuasi riak tegangan, disarankan memasang
kapasitor \textit{decoupling} $100\,\text{nF}$ secara paralel sedekat mungkin antara pin $V_{CC}$ dan GND.
Jika terdapat perbedaan level logika antara mikrokontroler (misal 5\,V) dengan komponen (3,3\,V),
jelaskan apakah modul telah memiliki sirkuit \textit{logic level converter} terintegrasi atau
memerlukan pembagi tegangan (\textit{voltage divider}) eksternal.

% ---------------------------------------------------------------------
\subsubsection{Spesifikasi Penting dari Datasheet}
Rangkum parameter-parameter penting yang diambil langsung dari datasheet pabrikan resmi:

\begin{table}[H]
	\centering
	\caption{Spesifikasi penting komponen Nama Komponen}
	\label{tab:nama_komponen_spesifikasi}
	\small
	\begin{tabularx}{\textwidth}{@{}Xcc@{}}
		\toprule
		Parameter                                           & Nilai                   & Satuan        \\
		\midrule
		Tegangan kerja ($V_{DD}$ / $V_{CC}$)                & $3{,}3 \text{--} 5{,}0$ & $\text{V DC}$ \\
		Arus konsumsi (Operasional Aktif)                   & $1{,}0 \text{--} 2{,}5$ & $\text{mA}$   \\
		Arus konsumsi (\textit{Standby / Sleep})            & $10 \text{--} 50$       & $\mu\text{A}$ \\
		Periode Pembaruan / Sampling Period                 & $1$                     & $\text{s}$    \\
		Suhu operasional lingkungan                         & $-40 \text{--} +85$     & $^\circ\text{C}$ \\
		Jarak Transmisi / Beban Maksimum                    & Hingga $20$             & $\text{m}$    \\
		\bottomrule
	\end{tabularx}
\end{table}

% ---------------------------------------------------------------------
\subsubsection{Koneksi dengan Mikrokontroler dan Implementasi Kode}
Jelaskan bagaimana komponen dihubungkan ke mikrokontroler target (seperti Arduino Mega 2560/ATmega2560
atau mikrokontroler lainnya). Sebutkan periferal internal mikrokontroler yang dialokasikan,
seperti pin GPIO digital, kanal ADC analog, pin PWM timer, interupsi eksternal, atau bus komunikasi I$^2$C/SPI.

\paragraph{Rangkaian Skematik dan Simulasi}
Pengujian antarmuka perangkat keras dapat divalidasi melalui simulator daring (seperti Wokwi)
maupun perakitan breadboard secara fisik:
\begin{itemize}
	\item \textbf{Koneksi Catu Daya:} Pin VCC dihubungkan ke rel daya $5\,\text{V}$ atau $3{,}3\,\text{V}$ MCU, dan GND dihubungkan ke ground sistem bersama (\textit{common ground}).
	\item \textbf{Koneksi Sinyal dan Komponen Pasif:} Pin sinyal dihubungkan ke pin I/O MCU yang sesuai. Jika menggunakan bus dengan saluran \textit{open-drain} (seperti 1-Wire atau I$^2$C), pasang resistor \textit{pull-up} eksternal bernilai $4{,}7\text{--}10\,\text{k}\Omega$ bila modul belum menyediakannya.
\end{itemize}

\begin{figure}[htbp]
	\centering
	% Ganti dengan file skematik rangkaian aktual atau tangkapan layar simulasi Wokwi
	% \includegraphics[width=0.70\textwidth]{images/skema_nama_komponen.png}
	\fbox{\parbox[c][5cm][c]{0.70\textwidth}{\centering Tempatkan tangkapan layar simulasi Wokwi atau diagram skematik di sini}}
	\caption{Rangkaian antarmuka Nama Komponen dengan mikrokontroler pada simulasi / perakitan fisik.}
	\label{fig:nama_komponen_schematic}
\end{figure}

Rincian pemetaan koneksi fisik antara pin komponen dan fungsi periferal mikrokontroler
dirangkum pada Tabel~\ref{tab:nama_komponen_mcu_mapping}.

\begin{table}[H]
	\centering
	\caption{Pemetaan antarmuka Nama Komponen ke mikrokontroler}
	\label{tab:nama_komponen_mcu_mapping}
	\small
	\begin{tabularx}{\textwidth}{@{}lXX@{}}
		\toprule
		Kebutuhan Komponen                & Fitur MCU                      & Alasan Pemilihan \& Fungsi                                   \\
		\midrule
		Jalur Transmisi Sinyal            & Pin GPIO / ADC / Timer PWM     & Membaca logika sinyal atau menghasilkan modulasi pemicu      \\
		Sinkronisasi Pewaktuan            & Hardware Timer / Delay System  & Menjaga presisi deteksi pulsa atau pewaktuan polling         \\
		Pencegahan Kondisi Mengambang     & Resistor Pull-Up Internal/Eksternal & Memastikan level tegangan mantap saat jalur komunikasi idle \\
		\bottomrule
	\end{tabularx}
\end{table}

\paragraph{Kode Demo Pengujian}
Berikut adalah contoh implementasi program minimal menggunakan pustaka resmi atau pembacaan
register langsung untuk membaca/mengendalikan komponen dan melaporkan hasilnya melalui komunikasi serial:

\begin{lstlisting}[language=C++, caption={Kode program pengujian Nama Komponen pada mikrokontroler}, label={lst:nama_komponen_code}]
// Header pustaka pendukung komponen
#define PIN_KOMPONEN 2 // Alokasi pin mikrokontroler

void setup() {
  Serial.begin(9600);
  Serial.println(F("Inisialisasi Pengujian Nama Komponen..."));

  // Konfigurasi pin atau inisialisasi pustaka
  pinMode(PIN_KOMPONEN, INPUT);
}

void loop() {
  // Interval pembacaan sesuai sampling rate komponen
  delay(1000);

  // Proses akuisisi data atau kendali komponen
  int nilaiSinyal = digitalRead(PIN_KOMPONEN);

  // Menampilkan hasil pembacaan ke Serial Monitor
  Serial.print(F("Status Komponen: "));
  Serial.println(nilaiSinyal);
}
\end{lstlisting}

\paragraph{Gambaran Hasil Pembacaan}
Setelah kode diunggah ke mikrokontroler atau dijalankan melalui simulasi, keluaran pembacaan data
pada \textit{Serial Monitor} akan tampak seperti pada contoh log berikut:

\begin{verbatim}
Inisialisasi Pengujian Nama Komponen...
Status Komponen: Aktif / Terbaca (Nilai = 24.50)
Status Komponen: Aktif / Terbaca (Nilai = 24.60)
Status Komponen: Aktif / Terbaca (Nilai = 24.55)
\end{verbatim}
